\documentclass[11pt]{article}
\usepackage{amsmath, amssymb}
\usepackage{proof}   % for \infer
\usepackage{geometry}
\usepackage{hyperref}
\usepackage{booktabs}

\geometry{margin=1in}

\newcommand{\typeof}[2]{#1 : #2}
\newcommand{\step}[2]{#1 \to #2}

\title{Lab 14: Evaluation and Typing Rules}
\author{}
\date{}

\begin{document}
\maketitle

\section*{Language Grammar}

\[
\begin{array}{rcll}
e & ::= & \texttt{true} \\
  & \mid & \texttt{false} \\
  & \mid & i & \text{(integer literals)} \\
  & \mid & s & \text{(string literals)} \\
  & \mid & \texttt{succ}\ e \\
  & \mid & \texttt{pred}\ e \\
  & \mid & \texttt{iszero}\ e \\
  & \mid & \texttt{if}\ e\ \texttt{then}\ e\ \texttt{else}\ e \\
  & \mid & \texttt{concat}\ e\ e \\
  & \mid & \texttt{isemptystr}\ e \\
  & \mid & \texttt{strlen}\ e \\[6pt]
T & ::= & \texttt{Bool} \mid \texttt{Int} \mid \texttt{String}
\end{array}
\]

\noindent\textbf{Values} (fully-evaluated, irreducible expressions):
\[
v\ ::=\ \texttt{true}\ \mid\ \texttt{false}\ \mid\ i\ \mid\ s
\]

% ────────────────────────────────────────────────
\section{Small-Step Evaluation Rules}
% ────────────────────────────────────────────────

\subsection*{\texttt{succ}}

\[
\infer[\textsc{E-Succ}]
  {\texttt{succ}\ e \to \texttt{succ}\ e'}
  {e \to e'}
\qquad
\infer[\textsc{E-SuccInt}]
  {\texttt{succ}\ i \to i+1}
  {}
\]

\subsection*{\texttt{pred}}

\[
\infer[\textsc{E-Pred}]
  {\texttt{pred}\ e \to \texttt{pred}\ e'}
  {e \to e'}
\qquad
\infer[\textsc{E-PredInt}]
  {\texttt{pred}\ i \to i-1}
  {}
\]

\subsection*{\texttt{iszero}}

\[
\infer[\textsc{E-IsZero}]
  {\texttt{iszero}\ e \to \texttt{iszero}\ e'}
  {e \to e'}
\qquad
\infer[\textsc{E-IsZeroTrue}]
  {\texttt{iszero}\ 0 \to \texttt{true}}
  {}
\qquad
\infer[\textsc{E-IsZeroFalse}\ (i \neq 0)]
  {\texttt{iszero}\ i \to \texttt{false}}
  {}
\]

\subsection*{\texttt{if}}

\[
\infer[\textsc{E-If}]
  {\texttt{if}\ e\ \texttt{then}\ e_1\ \texttt{else}\ e_2
   \to
   \texttt{if}\ e'\ \texttt{then}\ e_1\ \texttt{else}\ e_2}
  {e \to e'}
\]

\[
\infer[\textsc{E-IfTrue}]
  {\texttt{if}\ \texttt{true}\ \texttt{then}\ e_1\ \texttt{else}\ e_2 \to e_1}
  {}
\qquad
\infer[\textsc{E-IfFalse}]
  {\texttt{if}\ \texttt{false}\ \texttt{then}\ e_1\ \texttt{else}\ e_2 \to e_2}
  {}
\]

\subsection*{\texttt{concat}}

\[
\infer[\textsc{E-Concat1}]
  {\texttt{concat}\ e_1\ e_2 \to \texttt{concat}\ e_1'\ e_2}
  {e_1 \to e_1'}
\qquad
\infer[\textsc{E-Concat2}]
  {\texttt{concat}\ s\ e_2 \to \texttt{concat}\ s\ e_2'}
  {e_2 \to e_2'}
\]

\[
\infer[\textsc{E-ConcatStr}]
  {\texttt{concat}\ s_1\ s_2 \to s_1 \mathbin{+\!\!+} s_2}
  {}
\]

\noindent where $s_1 \mathbin{+\!\!+} s_2$ denotes string concatenation of literal values $s_1$ and $s_2$.

\subsection*{\texttt{isemptystr}}

\[
\infer[\textsc{E-IsEmptyStr}]
  {\texttt{isemptystr}\ e \to \texttt{isemptystr}\ e'}
  {e \to e'}
\]

\[
\infer[\textsc{E-IsEmptyStrTrue}]
  {\texttt{isemptystr}\ \text{``''} \to \texttt{true}}
  {}
\qquad
\infer[\textsc{E-IsEmptyStrFalse}\ (s \neq \text{``''})]
  {\texttt{isemptystr}\ s \to \texttt{false}}
  {}
\]

\subsection*{\texttt{strlen}}

\[
\infer[\textsc{E-StrLen}]
  {\texttt{strlen}\ e \to \texttt{strlen}\ e'}
  {e \to e'}
\qquad
\infer[\textsc{E-StrLenVal}]
  {\texttt{strlen}\ s \to |s|}
  {}
\]

\noindent where $|s|$ is the integer length of string literal $s$.

% ────────────────────────────────────────────────
\section{Typing Rules}
% ────────────────────────────────────────────────

\subsection*{Literals}

\[
\infer[\textsc{T-True}]
  {\Gamma \vdash \texttt{true} : \texttt{Bool}}
  {}
\qquad
\infer[\textsc{T-False}]
  {\Gamma \vdash \texttt{false} : \texttt{Bool}}
  {}
\qquad
\infer[\textsc{T-Int}]
  {\Gamma \vdash i : \texttt{Int}}
  {}
\qquad
\infer[\textsc{T-Str}]
  {\Gamma \vdash s : \texttt{String}}
  {}
\]

\subsection*{\texttt{succ} and \texttt{pred}}

\[
\infer[\textsc{T-Succ}]
  {\Gamma \vdash \texttt{succ}\ e : \texttt{Int}}
  {\Gamma \vdash e : \texttt{Int}}
\qquad
\infer[\textsc{T-Pred}]
  {\Gamma \vdash \texttt{pred}\ e : \texttt{Int}}
  {\Gamma \vdash e : \texttt{Int}}
\]

\subsection*{\texttt{iszero}}

\[
\infer[\textsc{T-IsZero}]
  {\Gamma \vdash \texttt{iszero}\ e : \texttt{Bool}}
  {\Gamma \vdash e : \texttt{Int}}
\]

\subsection*{\texttt{if}}

\[
\infer[\textsc{T-If}]
  {\Gamma \vdash \texttt{if}\ e\ \texttt{then}\ e_1\ \texttt{else}\ e_2 : T}
  {\Gamma \vdash e : \texttt{Bool}
   & \Gamma \vdash e_1 : T
   & \Gamma \vdash e_2 : T}
\]

\subsection*{\texttt{concat}}

\[
\infer[\textsc{T-Concat}]
  {\Gamma \vdash \texttt{concat}\ e_1\ e_2 : \texttt{String}}
  {\Gamma \vdash e_1 : \texttt{String}
   & \Gamma \vdash e_2 : \texttt{String}}
\]

\subsection*{\texttt{isemptystr}}

\[
\infer[\textsc{T-IsEmptyStr}]
  {\Gamma \vdash \texttt{isemptystr}\ e : \texttt{Bool}}
  {\Gamma \vdash e : \texttt{String}}
\]

\subsection*{\texttt{strlen}}

\[
\infer[\textsc{T-StrLen}]
  {\Gamma \vdash \texttt{strlen}\ e : \texttt{Int}}
  {\Gamma \vdash e : \texttt{String}}
\]

\section*{Progress and Preservation}

\textbf{Progress.}
If $\typeof{e}{T}$, then $e$ is a value or $\exists\, e'$ with $\step{e}{e'}$.

The typing rules enforce that every operator receives the right kind of argument.
Because no well-typed expression can be an irreducible non-value, evaluation always has somewhere to go:
integer operators (\texttt{succ}, \texttt{pred}, \texttt{iszero}) are only ever applied to integers,
string operators (\texttt{concat}, \texttt{isemptystr}, \texttt{strlen}) only to strings,
and the condition of \texttt{if} is always a boolean.

\medskip
\textbf{Preservation.}
If $\typeof{e}{T}$ and $\step{e}{e'}$, then $\typeof{e'}{T}$.

The type of an expression is invariant under evaluation because the computation rules produce
results of the same type that the typing rules predict:
integer arithmetic stays \texttt{Int}, comparisons and boolean tests produce \texttt{Bool},
string concatenation stays \texttt{String}, \texttt{strlen} yields \texttt{Int},
and both branches of \texttt{if} are required to share the same type $T$,
so whichever branch is chosen, the result still has type $T$.

\end{document}
